% Глава: 20-dockerfile
% Статус: наполнено (первый прогон).
\chapter{Dockerfile}
\label{ch:20-dockerfile}

\section{Что такое Dockerfile}

\textbf{Dockerfile} --- текст с инструкциями, как
построить образ. Команда \cmd{docker build} выполняет
инструкции по порядку, каждая --- отдельный
\textbf{слой}.

Собираем образы на своей ВМ: \cmd{docker build} запускается на
\fpath{pglab.nsu.ru} и кэширует слои на её диске.

\begin{minted}[bgcolor=gray!6]{Dockerfile}
# базовый образ
FROM ubuntu:22.04

# метаданные
LABEL maintainer="dev@example.org"

# рабочий каталог внутри контейнера
WORKDIR /app

# установить инструменты; каждый RUN создаёт отдельный слой
RUN apt-get update && apt-get install -y --no-install-recommends gcc make
RUN rm -rf /var/lib/apt/lists/*

# скопировать файлы проекта
COPY src/ ./src/

# собрать программу
RUN make

# точка входа
CMD ["./app"]
\end{minted}

\section{Основные инструкции}

\begin{center}
\begin{tabular}{lp{0.62\linewidth}}
\toprule
\textbf{Инструкция} & \textbf{Назначение} \\
\midrule
\cmd{FROM}     & базовый образ; начинается с него сборка \\
\cmd{WORKDIR}  & каталог по умолчанию внутри контейнера \\
\cmd{RUN}      & выполнить команду при сборке (создаёт слой) \\
\cmd{COPY}     & скопировать файлы из контекста сборки \\
\cmd{ADD}      & как COPY + распаковка архивов и загрузка по URL \\
\cmd{ENV}      & задать переменную окружения в образе \\
\cmd{ARG}      & параметр, передаваемый при \cmd{docker build} \\
\cmd{LABEL}    & метаданные образа \\
\cmd{EXPOSE}   & документация о портах (не пробрасывает!) \\
\cmd{ENTRYPOINT} & фиксированная команда контейнера \\
\cmd{CMD}      & аргументы по умолчанию для ENTRYPOINT \\
\cmd{USER}     & пользователь, от которого работает контейнер \\
\bottomrule
\end{tabular}
\end{center}

\textbf{CMD} vs \textbf{ENTRYPOINT}: \cmd{CMD} легко
переопределить аргументами запуска; \cmd{ENTRYPOINT} ---
``главная'' команда, аргументы передаются ей.

\section{Слои и кэш}

Каждая инструкция создаёт слой. Docker кэширует
слои: если строки выше не менялись, их слои берутся
из кэша, и сборка идёт быстро. Поэтому
\textbf{порядок инструкций важен}:
\begin{itemize}
  \item сначала то, что редко меняется (зависимости);
  \item потом то, что часто (исходный код);
  \item так при изменении кода пересобирается только
    верхний слой.
\end{itemize}

\section{\cmd{COPY} и \cmd{ADD}}

\cmd{COPY} --- скопировать файл из каталога сборки.
\cmd{ADD} делает то же, но дополнительно распаковывает
архивы и может загрузить файл по URL. Для обычного
копаирования файлов проекта используйте \cmd{COPY}.

\section{Многоступенчатая сборка (multi-stage)}

Один Dockerfile может содержать несколько
\cmd{FROM}:
\begin{minted}[bgcolor=gray!6]{Dockerfile}
FROM golang:1.21 AS build
WORKDIR /src
COPY . .
RUN go build -o app .

FROM alpine:3.18
COPY --from=build /src/app /app
ENTRYPOINT ["/app"]
\end{minted}

Первый слой --- «сборочный» (с компилятором), второй ---
лёгкий, только с готовым бинарником. Размер
конечного образа резко меньше.

\section{\cmd{.dockerignore}}

Файл \fpath{.dockerignore} рядом с Dockerfile
определяет, что \emph{не} копировать в контекст
сборки (.git, локальные логи, артефакты).
Уменьшает размер и время сборки.

\section{Типичные ошибки}
\begin{itemize}
  \item \cmd{RUN apt-get update} и \cmd{apt-get install}
    в разных \cmd{RUN} --- кэш не поможет; объедините
    в одну инструкцию;
  \item Забыли \cmd{WORKDIR} --- файлы копировались в
    \fpath{/} вместо нужного каталога;
  \item Образ «огромный» --- кэш пакетного менеджера не
    чистится после установки (см.\ пример с \cmd{rm -rf}
    выше), либо не используется multi-stage;
  \item \cmd{EXPOSE} «не работает» --- он не пробрасует
    порт; нужен \cmd{-p} при \cmd{docker run}.
\end{itemize}

\section{Связь с остальным курсом}

Dockerfile --- описание образа; хранение данных ---
\nameref{ch:21-docker-data}; несколько контейнеров ---
\nameref{ch:22-docker-compose}; сборка C-проекта внутри
образа --- \nameref{ch:13-c-compilation} и
\nameref{ch:14-make}.

\section{Задачи для самопроверки}
\begin{enumerate}
  \item Чем \cmd{COPY} отличается от \cmd{ADD}?
  \item Что такое слой и почему порядок инструкций
    важен?
  \item Чем \cmd{CMD} отличается от \cmd{ENTRYPOINT}?
  \item Как сделать образ «меньше», используя
    multi-stage?
  \item Зачем нужен \cmd{.dockerignore}?
\end{enumerate}
